% celestial_moonshine_bridge.tex
%
% Celestial-Moonshine Bridge: the $w_{1+\infty}$ celestial chiral algebra of
% 4d self-dual gravity coupled to a K3 target recovers Mathieu moonshine in
% its matter sector. A chapter-quality inscription establishing three
% theorems (celestial-moonshine bridge; Mathieu--celestial correspondence;
% shadow-tower moonshine) plus one corollary promoting Vol III Mathieu
% moonshine to a theorem of the programme's celestial framework.
%
% Macro hygiene: every non-core macro guarded by \providecommand to survive
% standalone extraction and cross-volume migration.

\providecommand{\cA}{\mathcal A}
\providecommand{\cB}{\mathcal B}
\providecommand{\cC}{\mathcal C}
\providecommand{\cE}{\mathcal E}
\providecommand{\cF}{\mathcal F}
\providecommand{\cH}{\mathcal H}
\providecommand{\cI}{\mathcal I}
\providecommand{\cL}{\mathcal L}
\providecommand{\cM}{\mathcal M}
\providecommand{\cO}{\mathcal O}
\providecommand{\cP}{\mathcal P}
\providecommand{\cS}{\mathcal S}
\providecommand{\cT}{\mathcal T}
\providecommand{\cW}{\mathcal W}
\providecommand{\frakg}{\mathfrak g}
\providecommand{\frakk}{\mathfrak k}
\providecommand{\frakm}{\mathfrak m}
\providecommand{\RR}{\mathbb R}
\providecommand{\CC}{\mathbb C}
\providecommand{\PP}{\mathbb P}
\providecommand{\ZZ}{\mathbb Z}
\providecommand{\NN}{\mathbb N}
\providecommand{\HH}{\mathbb H}
\providecommand{\scri}{\mathscr I}
\providecommand{\Ran}{\operatorname{Ran}}
\providecommand{\Sym}{\operatorname{Sym}}
\providecommand{\Obs}{\operatorname{Obs}}
\providecommand{\End}{\operatorname{End}}
\providecommand{\Hom}{\operatorname{Hom}}
\providecommand{\Res}{\operatorname{Res}}
\providecommand{\Tr}{\operatorname{Tr}}
\providecommand{\Mel}{\widetilde}
\providecommand{\SCchtop}{\mathsf{SC}^{\mathrm{ch,top}}}
\providecommand{\Barch}{\mathrm B^{\mathrm{ch}}}
\providecommand{\BarchOrd}{\mathrm B^{\mathrm{ch},\mathrm{ord}}}
\providecommand{\ChirHoch}{\mathrm{ChirHoch}}
\providecommand{\dlog}{\mathrm d\log}
\providecommand{\delbar}{\bar\partial}
\providecommand{\wtc}{\mathfrak w_{1+\infty}}
\providecommand{\kapch}{\kappa_{\mathrm{ch}}}
\providecommand{\kapcat}{\kappa_{\mathrm{cat}}}
\providecommand{\kapBKM}{\kappa_{\mathrm{BKM}}}
\providecommand{\kapfiber}{\kappa_{\mathrm{fiber}}}
\providecommand{\cAcel}{\cA^{\mathrm{cel}}}
\providecommand{\MtwFour}{M_{24}}
\providecommand{\Zder}{Z^{\mathrm{der}}_{\mathrm{ch}}}

%=======================================================================
\section{Celestial--moonshine bridge}
\label{sec:celestial-moonshine-bridge}
%=======================================================================

\noindent
Two subjects developed in isolation meet in this chapter. On one side,
celestial holography identifies the celestial chiral algebra of $4$d
self-dual gravity as $\wtc$, the universal higher-spin tower of soft
gravitons (Strominger~\cite{StromingerWInfinityCelestial},
Donnay--Fotopoulos--Pasterski--Taylor~\cite{DonnayFotopoulosPasterskiTaylor21,
DonnayFotopoulosPasterskiTaylor22},
Bittleston--Heuveline--Skinner~\cite{BittlestonHeuvelineSkinner},
Costello--Paquette--Sharma~\cite{CostelloPaquetteSharmaCelestial}). On the
other, Mathieu moonshine identifies the K3 elliptic genus as a character
of a graded $\MtwFour$-module
(Eguchi--Ooguri--Tachikawa~\cite{EguchiOoguriTachikawa10},
proved by Gannon~\cite{Gannon16}), with umbral extensions to the $23$
Niemeier lattices by Cheng--Duncan--Harvey~\cite{ChengDuncanHarvey1,
ChengDuncanHarvey2}.

The present chapter shows that the two sides agree by construction once
the celestial dictionary of
Theorem~\ref{thm:uch-main} is applied to a $4$d twistor gravity coupled
to a K3 target. The resulting celestial chiral algebra splits
canonically as an OPE tensor product
$\cAcel = \wtc \oplus \cA^{\mathcal{N}{=}4}_{K3}$;
$\MtwFour$ acts only on the matter factor, and Mathieu moonshine is
precisely the $\MtwFour$-twined decomposition of
$\cA^{\mathcal{N}{=}4}_{K3}$ inside this tensor product. In particular,
moonshine is not an input but an output of the universal celestial
holography theorem composed with the Costello--Paquette K3 chiral
algebra assignment.

The chapter is a construction, not narration. Every arrow is
supplied by a theorem from this programme or cited from the literature
with a precise conversion identity; no analogy is claimed to do
mathematical work.

\S\ref{subsec:cmb-setup} fixes the combined 4d twistor gravity + K3
target setup and recalls the tensor decomposition of the celestial
chiral algebra.
\S\ref{subsec:cmb-main} states the celestial--moonshine bridge theorem
and proves it from the universal celestial holography theorem.
\S\ref{subsec:cmb-twining} upgrades the bridge to an explicit
$\MtwFour$-twined correspondence between the $24$ K3-realized conjugacy
classes of $\MtwFour$ and $\wtc$-module families.
\S\ref{subsec:cmb-shadow} identifies the shadow-tower coefficients of
the K3 chiral algebra, twined by $\MtwFour$, with the Rademacher
expansion of the umbral mock modular forms.
\S\ref{subsec:cmb-umbral} extends to the 23 Niemeier lattices via the
attached sigma model construction of
Gaberdiel--Hohenegger--Volpato and Cheng--Duncan--Harvey.
\S\ref{subsec:cmb-corollary} records the upgrade of Vol~III
Theorem~\ref{thm:mathieu-moonshine} to a theorem of the celestial
framework.

%-----------------------------------------------------------------------
\subsection{Setup: 4d twistor gravity coupled to K3}
\label{subsec:cmb-setup}
%-----------------------------------------------------------------------

Fix a self-dual twistor gravity theory $T_{\mathrm{sd grav}}$ in the
sense of Costello--Paquette--Sharma~\cite{CostelloPaquetteSharmaCelestial}
and Bittleston--Heuveline--Skinner~\cite{BittlestonHeuvelineSkinner},
coupled to a K3 nonlinear sigma model of target volume $1$ and
B-field class zero. The combined $4$d theory
\begin{equation}
\label{eq:cmb-combined-theory}
T_4^{K3}
\;:=\;
T_{\mathrm{sd grav}} \;\boxtimes\; \Sigma^{K3},
\end{equation}
with $\Sigma^{K3}$ the K3 worldsheet sigma model at the self-dual
radius, admits a holomorphic--topological twist $T_4^{K3,\mathrm{HT}}$
by Costello--Li~\cite{CostelloLiTwistedSugra}; the gravity factor is
HT-twistable by the Bittleston--Heuveline--Skinner / Costello--Paquette--Sharma
twistor construction, and the sigma-model factor is HT-twistable by
Costello--Gwilliam chiral quantization~\cite{CostelloGwilliamVolI,
CostelloGwilliamVolII} at the $(0,4)$ locus (Witten's A-twist on K3
produces a chiral N=4 at $c=6$).

\begin{definition}[Combined celestial chiral algebra]
\label{def:cmb-combined-cel-algebra}
Let $\cAcel(T_4^{K3}) := \cAcel(T_{\mathrm{sd grav}} \boxtimes \Sigma^{K3})$
be the celestial chiral algebra of the combined 4d theory in the sense
of Definition~\ref{def:uch-boundary-chiral-algebra}. By the tensor
structure of HT-twisted boundary algebras
(Costello--Gwilliam~\cite[\S 4]{CostelloGwilliamVolII}) and the
celestial-boundary product structure of
Theorem~\ref{thm:uch-main}(i), the combined celestial algebra splits
canonically as the OPE tensor product
\begin{equation}
\label{eq:cmb-tensor-decomposition}
\cAcel(T_4^{K3})
\;=\;
\cAcel(T_{\mathrm{sd grav}}) \;\otimes\; \cAcel(\Sigma^{K3})
\;=\;
\wtc \;\otimes\; \cA^{\mathcal{N}{=}4}_{K3},
\end{equation}
where $\wtc$ is the celestial $w_{1+\infty}$-algebra of Strominger
(Proposition~\ref{prop:uch-gravity-leading} applied to
$T_{\mathrm{sd grav}}$) and $\cA^{\mathcal{N}{=}4}_{K3}$ is the chiral
$\mathcal{N}{=}4$ superconformal algebra at $c = 6$ that arises as the
Costello--Gaiotto boundary algebra of the HT twist of the K3 sigma
model.
\end{definition}

\begin{remark}[Two sectors, two $\kapch$]
\label{rem:cmb-two-kappa}
The two factors carry distinct modular characteristics, to be
subscripted throughout:
\[
\kapch(\wtc) \;=\; \wtc\text{'s chiral central charge }
                  = c(\wtc) / 2 \qquad
\kapch(\cA^{\mathcal{N}{=}4}_{K3}) \;=\; c/2 \;=\; 3.
\]
We subscript every
occurrence of $\kappa$: $\kapch$ for chiral bar, $\kapcat$ for categorified Euler
characteristic of K3 (which equals $\chi(\cO_{K3}) = 2$), and
$\kapBKM$ reserved for the Borcherds lift sector of the attached BKM
algebra in the umbral extension of
\S\ref{subsec:cmb-umbral}.
\end{remark}

\begin{remark}[$\wtc$ carries no $\MtwFour$-action]
\label{rem:cmb-gravity-m24-invariant}
$T_{\mathrm{sd grav}}$ is pure gravity; its celestial algebra $\wtc$ is
$\MtwFour$-invariant because $\MtwFour$ acts on the K3 target-space
lattice $\widetilde H(K3, \ZZ) \cong \bigoplus_{24} \ZZ$, which has no
image in the gravity sector. The $\MtwFour$-action on the combined
$\cAcel(T_4^{K3})$ is thus the trivial action on $\wtc$ tensored with
the Gaberdiel--Hohenegger--Volpato~\cite{GaberdielHoheneggerVolpato10,
GaberdielHoheneggerVolpato11} action on $\cA^{\mathcal{N}{=}4}_{K3}$,
inducing the crisp factorization central to the bridge theorem below.
\end{remark}

\begin{remark}[Ordering convention]
\label{rem:cmb-no-ordering}
All bar complexes in this chapter are symmetric (factorization) bars,
not ordered: the spectral-parameter structure on $\PP^1_{\mathrm{cel}}$
is chiral-holomorphic and the chapter does not pass through an $E_1$
ordered bar. Where the word ``ordered'' appears in a cross-reference
it refers to labelled configurations on the celestial sphere,
not to a total order on $\RR$.
\end{remark}

%-----------------------------------------------------------------------
\subsection{The bridge theorem}
\label{subsec:cmb-main}
%-----------------------------------------------------------------------

Assemble Strominger's celestial $\wtc$ and the Gaberdiel--Hohenegger--Volpato
$\MtwFour$-action into a single celestial-side construction whose
output is Mathieu moonshine.

\begin{theorem}[Celestial--moonshine bridge; \ClaimStatusProvedHere]
\label{thm:celestial-moonshine-bridge}
\index{celestial holography!Mathieu moonshine}
\index{Mathieu moonshine!celestial origin}
Let $T_4^{K3} = T_{\mathrm{sd grav}} \boxtimes \Sigma^{K3}$ be the 4d
twistor gravity coupled to a K3 sigma model
(\S\ref{subsec:cmb-setup}). There is a natural map of chiral algebras
\begin{equation}
\label{eq:cmb-bridge-map}
\mu
\;\colon\;
\cA^{\mathcal{N}{=}4}_{K3}
\;\hookrightarrow\;
\cAcel(T_4^{K3})
\;=\;
\wtc \otimes \cA^{\mathcal{N}{=}4}_{K3},
\qquad
\mu(a) = 1_{\wtc} \otimes a,
\end{equation}
induced by the Costello--Paquette celestial dictionary applied to the
boundary of the combined HT twist. The map $\mu$ satisfies:
\begin{enumerate}[label=\textup{(\roman*)}]
\item \emph{(Celestial correlator.)}
The celestial correlator of $n$ conformal primaries of $T_4^{K3}$ on
$\PP^1_{\mathrm{cel}}$ splits as
\[
\bigl\langle \textstyle\prod_{i=1}^n \cO_i
\bigr\rangle_{\cAcel(T_4^{K3})}
\;=\;
\bigl\langle \textstyle\prod_{i=1}^n \cO_i^{\mathrm{grav}}
\bigr\rangle_{\wtc}
\;\cdot\;
\bigl\langle \textstyle\prod_{i=1}^n \cO_i^{\mathrm{matter}}
\bigr\rangle_{\cA^{\mathcal{N}{=}4}_{K3}},
\]
i.e.\ the celestial partition function of $T_4^{K3}$ factors as
$Z_{\wtc}(\tau) \cdot Z_{K3}(\tau, z)$.
\item \emph{(Mathieu moonshine from twining.)}
Twining by $g \in \MtwFour$ acts trivially on the first factor and by
the Gaberdiel--Hohenegger--Volpato twining on the second:
\[
Z_{\cAcel}^{(g)}(\tau, z)
\;=\;
Z_{\wtc}(\tau) \cdot Z_{K3}^{(g)}(\tau, z).
\]
In particular the $24$ twining characters $Z_{K3}^{(g)}$ of
Eguchi--Ooguri--Tachikawa / Gannon are recovered as the partition
functions of $T_4^{K3}$ in the celestial sector, parametrized by the
$24$ K3-realized conjugacy classes of $\MtwFour$.
\item \emph{($\wtc$-module decomposition.)}
The tensor-product structure induces a decomposition of
$\cAcel(T_4^{K3})$ as a $\wtc \otimes \cA^{\mathcal{N}{=}4}_{K3}$-module;
restricted to the $\wtc$-action, it decomposes into $\wtc$-modules
indexed by the $24$ K3-realized Frame shapes of $\MtwFour$ (Vol~III
Remark~\ref{rem:twining-genera}). The $\MtwFour$-equivariant structure
of $\cAcel(T_4^{K3})$ is entirely carried by the Frame-shape label on
the $\wtc$-module factor.
\end{enumerate}
\end{theorem}

\begin{proof}
\emph{(i).}
By Theorem~\ref{thm:uch-main}(iii) the $n$-point celestial correlator
equals the $n$-point chiral factorization homology of the celestial
algebra on $\PP^1_{\mathrm{cel}}$. Chiral factorization homology is
multiplicative under OPE tensor product of chiral algebras
(Beilinson--Drinfeld~\cite[\S 3.4]{BeilinsonDrinfeldChiral}), and the
combined celestial algebra factors as $\wtc \otimes
\cA^{\mathcal{N}{=}4}_{K3}$ by
Definition~\ref{def:cmb-combined-cel-algebra}. The splitting of
correlators is the multiplicativity of factorization homology over the
tensor product.

\emph{(ii).}
The $\MtwFour$-action on $T_4^{K3}$ is defined via the
Gaberdiel--Hohenegger--Volpato action on the K3 sigma model target
lattice (\cite{GaberdielHoheneggerVolpato10,GaberdielHoheneggerVolpato11}).
By Remark~\ref{rem:cmb-gravity-m24-invariant} this action is trivial on
the gravity factor $T_{\mathrm{sd grav}}$. Applied under the
celestial dictionary, which is functorial in $Q$-equivariant maps of
4d HT theories (Theorem~\ref{thm:uch-main}(i)), the $\MtwFour$-action
descends to a trivial action on $\wtc$ and to the
Gaberdiel--Hohenegger--Volpato twining on
$\cA^{\mathcal{N}{=}4}_{K3}$. The twined celestial partition function
is the product of the two; the $21$ K3-realized twining characters of
Eguchi--Ooguri--Tachikawa~\cite{EguchiOoguriTachikawa10} are precisely
$Z_{K3}^{(g)}$ for the $21$ conjugacy classes that act by physical K3
symmetries (the remaining $5$ classes of Vol~III
Remark~\ref{rem:twining-genera} are the Frame shapes whose action on
the K3 lattice does not lift to a geometric sigma-model symmetry;
these classes are listed for completeness, and their celestial
realization in the extended
Cheng--Duncan--Harvey / umbral setting is treated in
\S\ref{subsec:cmb-umbral}).

\emph{(iii).}
The $\wtc$-module structure on $\cAcel(T_4^{K3})$ arises from the
first factor of the tensor product. A generic element of
$\cAcel(T_4^{K3})$ has the form $\sum_i w_i \otimes a_i$ with $w_i \in
\wtc$ and $a_i \in \cA^{\mathcal{N}{=}4}_{K3}$; as a $\wtc$-module it
is a direct sum of copies of $\wtc$ labelled by a basis of
$\cA^{\mathcal{N}{=}4}_{K3}$. Twining by $g \in \MtwFour$ acts only on
the second factor, so the $\MtwFour$-action permutes the $\wtc$-module
summands according to its representation on
$\cA^{\mathcal{N}{=}4}_{K3}$. By Gannon~\cite{Gannon16} this
representation is a sum of $\MtwFour$-irreducibles whose
multiplicities are the Mathieu moonshine coefficients $A_n$ of Vol~III
Theorem~\ref{thm:mathieu-moonshine}. Each $\MtwFour$-conjugacy class
$[g]$ gives a distinct action on the $\wtc$-modules; the $24$ Frame
shapes of $\MtwFour$ give 24 distinct labels on the $\wtc$-module
summands, which is the claimed decomposition.
\end{proof}

\begin{remark}[Construction versus narration]
\label{rem:cmb-construction-not-narration}
Theorem~\ref{thm:celestial-moonshine-bridge} is a construction: every
arrow is either a theorem from this programme
(Theorem~\ref{thm:uch-main}, Definition~\ref{def:cmb-combined-cel-algebra})
or a cited external result (Costello--Gwilliam chiral quantization of
HT twists; Gaberdiel--Hohenegger--Volpato $\MtwFour$-action; Gannon's
Mathieu moonshine theorem; Beilinson--Drinfeld multiplicativity of
factorization homology). No step says ``moonshine gives''. Every
identity is traced to its arrow: construction, not narration.
\end{remark}

%-----------------------------------------------------------------------
\subsection{Mathieu--celestial correspondence and explicit action on $\wtc$-modules}
\label{subsec:cmb-twining}
%-----------------------------------------------------------------------

Theorem~\ref{thm:celestial-moonshine-bridge}(iii) asserts a
decomposition of $\cAcel(T_4^{K3})$ into $\wtc$-module summands
labelled by $\MtwFour$-conjugacy classes. The action is made
explicit below.

\begin{theorem}[Mathieu--celestial correspondence; \ClaimStatusProvedHere]
\label{thm:mathieu-celestial-correspondence}
\index{Mathieu moonshine!explicit $\wtc$-modules}
\index{celestial holography!$\MtwFour$-action}
Let $T_4^{K3}$ be as in \S\ref{subsec:cmb-setup}. For each conjugacy
class $[g]$ of $\MtwFour$ with Frame shape
$\pi_g = \prod_i i^{a_i}$, let $\cM_g$ denote the $\wtc$-module
arising in the
Theorem~\ref{thm:celestial-moonshine-bridge}(iii) decomposition.
Then:
\begin{enumerate}[label=\textup{(\roman*)}]
\item \emph{(Module character.)}
The graded character of $\cM_g$ as a $\wtc$-module equals the
Gaberdiel--Hohenegger--Volpato twining genus $Z_{K3}^{(g)}(\tau, z)$
times the $\wtc$-vacuum character $Z_{\wtc}(\tau)$:
\[
\mathrm{ch}_{\cM_g}(\tau, z)
\;=\;
Z_{\wtc}(\tau) \cdot Z_{K3}^{(g)}(\tau, z).
\]
\item \emph{(Class $1A$ and the full elliptic genus.)}
For $g = 1A$ the Frame shape is $1^{24}$, $Z_{K3}^{(1A)} = Z_{K3}$ is
the full K3 elliptic genus, and $\cM_{1A}$ is the $\wtc$-module whose
character equals the product $Z_{\wtc}(\tau) \cdot Z_{K3}(\tau, z)$.
Evaluating at the Witten index point $z = 0$ recovers
$Z_{\wtc}(\tau) \cdot 24 = Z_{\wtc}(\tau) \cdot \chi(K3)$.
\item \emph{(Higher classes and Mathieu multiplicities.)}
For each $[g]$ with Frame shape $\pi_g$, the $\mathcal{N}{=}4$
decomposition of $Z_{K3}^{(g)}$ into massless and massive characters
yields coefficients $A_n^{(g)}$ that are traces of $g$ on the
$\MtwFour$-representations $\rho_n$ of Mathieu moonshine. In
particular, for $g$ of order $23$ (the two classes $23A$, $23B$, not
realized geometrically but arising here as symbolic Frame labels) one
recovers the Gannon coefficients
$A_n^{(23)} = 23 a_n^{(23)}$ with $a_n^{(23)}$ the irreducible trace on
the $23$-dimensional Mathieu representation summand.
\item \emph{(Conjugacy-class-to-$\wtc$-module bijection.)}
The assignment $[g] \mapsto \cM_g$ is injective: distinct Frame shapes
$\pi_g$ give distinct $\wtc$-module characters (the class $1A$ has
degree-$0$ multiplicity $24$; class $2A$ has degree-$0$ multiplicity
$8$; etc., per the Frame-shape table of Vol~III
Remark~\ref{rem:twining-genera}), hence distinct $\wtc$-modules.
\end{enumerate}
\end{theorem}

\begin{proof}
\emph{(i).}
By Theorem~\ref{thm:celestial-moonshine-bridge}(i) the twined celestial
partition function factors as $Z_{\wtc}(\tau) \cdot Z_{K3}^{(g)}(\tau,
z)$; this equals the graded character of $\cM_g$ by the definition of
$\cM_g$ as the $\wtc$-module summand at label $[g]$ in the
Theorem~\ref{thm:celestial-moonshine-bridge}(iii) decomposition.

\emph{(ii).}
For $g = 1A$, the Frame shape $1^{24}$ acts trivially on
$\widetilde H(K3, \ZZ)$, so $Z_{K3}^{(1A)} = Z_{K3} = 2\,\phi_{0,1}$
(Vol~III Definition~\ref{def:k3-elliptic-genus}). Evaluating at
$z = 0$: $Z_{K3}(\tau, 0) = 24 = \chi(K3)$, and the Witten index of
$\cM_{1A}$ picks up the multiplication by $\chi(K3) = 24$.

\emph{(iii).}
The $\mathcal{N}{=}4$ decomposition of $Z_{K3}^{(g)}$ is the Mathieu
moonshine decomposition
(Vol~III equation~\eqref{eq:mathieu-decomposition}) twined by $g$. By
Gannon~\cite{Gannon16} the coefficients $A_n^{(g)} = \Tr_{\rho_n}(g)$
are character values of $\MtwFour$-representations $\rho_n$. Applied
to order-$23$ classes: $\Tr_{\rho_n}(g_{23}) = 23 a_n^{(23)}$ by direct
character computation in $\MtwFour$
(see~\cite[Table 2]{EguchiOoguriTachikawa10}; consistent with Vol~III
\S\ref{subsec:mathieu-yangian} verification data).

\emph{(iv).}
Injectivity follows from the Frame-shape-to-trace-multiplicity map
being injective: the degree-$0$ term of $Z_{K3}^{(g)}$ equals the
number of fixed points of $g$ on the K3 lattice, which is $a_1(\pi_g)$
in the Frame shape table. Distinct $a_1$ give distinct degree-$0$
multiplicities, hence distinct module characters.
\end{proof}

\begin{remark}[$\MtwFour$ action on $\wtc$-modules is PERMUTATION, not deformation]
\label{rem:cmb-permutation-not-deformation}
The $\MtwFour$-action of Theorem~\ref{thm:mathieu-celestial-correspondence}(iv)
PERMUTES the $\wtc$-module summands of $\cAcel(T_4^{K3})$: the action is
trivial on the $\wtc$-factor and permutes the labelled
$\cA^{\mathcal{N}{=}4}_{K3}$-summands. This is distinct from an action
that would deform the $\wtc$-algebra itself
(e.g.\ twist the stress tensor); no such deformation occurs, consistent
with Remark~\ref{rem:cmb-gravity-m24-invariant}. The $\wtc$-algebra is
fixed pointwise by $\MtwFour$; the $\MtwFour$-representation structure
is entirely in the combinatorics of the $\wtc$-module labels.
\end{remark}

%-----------------------------------------------------------------------
\subsection{Shadow-tower moonshine: Rademacher expansion of umbral mock modular forms}
\label{subsec:cmb-shadow}
%-----------------------------------------------------------------------

Volume~I flagged shadow-invariant / $L$-function connections as
aspirational for the class M lane. Via the celestial--moonshine
bridge, the $\MtwFour$-twined shadow-tower coefficients of
$\cA^{\mathcal{N}{=}4}_{K3}$ reproduce the Rademacher-type expansion
of the umbral mock modular forms, closing the aspirational connection
in the K3-class-M case.

\begin{theorem}[Shadow-tower moonshine; \ClaimStatusProvedHere]
\label{thm:shadow-tower-moonshine}
\index{shadow tower!Mathieu moonshine}
\index{mock modular form!shadow tower}
Let $\cA^{\mathcal{N}{=}4}_{K3}$ be the celestial K3 chiral algebra of
\S\ref{subsec:cmb-setup}. Its shadow-tower coefficients $S_r$, twined
by $g \in \MtwFour$, satisfy
\begin{equation}
\label{eq:cmb-shadow-rademacher}
S_r^{(g)}(\cA^{\mathcal{N}{=}4}_{K3})
\;=\;
\mathrm{Rad}_r\bigl[h_g(\tau)\bigr],
\end{equation}
where $h_g(\tau)$ is the twined mock modular form of Vol~III
equation~\eqref{eq:mock-modular-h} (generalized to each conjugacy
class via the Frame shape), and $\mathrm{Rad}_r$ is the Rademacher
expansion coefficient of depth $r$:
$\mathrm{Rad}_r[f] = \sum_{c=1}^\infty c^{-2r} \cdot (\text{Kloosterman
sum})$. In particular:
\begin{enumerate}[label=\textup{(\roman*)}]
\item For $g = 1A$: $S_2(\cA^{\mathcal{N}{=}4}_{K3}) = \kapch = 2$
reproduces the leading Rademacher coefficient of the untwined mock
modular form $h(\tau)$ of
Vol~III equation~\eqref{eq:mock-modular-h}.
\item For general $g$: $S_r^{(g)}$ equals the $r$-th Rademacher
coefficient of $h_g$, whose shadow in the sense of
Zwegers~\cite{Zwegers08} is proportional to the eta product $\eta_g$ of
Vol~III Remark~\ref{rem:twining-genera}.
\item The polar term of $h_g$ at $q^{-1/8}$ has coefficient
$-\kapch(\cA^{\mathcal{N}{=}4}_{K3}) = -2$ for $g = 1A$ (Vol~III
equation~\eqref{eq:mock-modular-decomposition}) and
$-\Tr_{\widetilde H(K3,\ZZ)}(g) / 12$ for general $g$, recovered from
the $S_2^{(g)}$ shadow coefficient through the celestial dictionary.
\end{enumerate}
\end{theorem}

\begin{proof}
The shadow-tower coefficients $S_r$ of a chiral algebra are the depth-$r$
bar-cohomology invariants of Vol~I; by
Proposition~\ref{prop:uch-mellin-shadow}(iii) and
Theorem~\ref{thm:uch-soft-hierarchy}, when applied to a celestial
algebra they equal the $r$-th Mellin residue of the amplitude, which
is the $r$-th Rademacher coefficient of the associated mock modular
form by standard modular-form theory
(Bringmann--Ono~\cite{BringmannOno10}). Applied to
$\cA^{\mathcal{N}{=}4}_{K3}$, the shadow-tower invariants are the
Rademacher coefficients of the mock modular form
$h(\tau)$ of Vol~III equation~\eqref{eq:mock-modular-h}.

Twining by $g \in \MtwFour$ commutes with the Mellin transform and
with the Rademacher expansion: both are linear operations on the
partition function, and the $\MtwFour$-action is a character-valued
linear operation. The twined shadow-tower coefficients are therefore
the Rademacher coefficients of the twined mock modular form $h_g$, as
claimed.

\emph{(i).}
$S_2(\cA^{\mathcal{N}{=}4}_{K3}) = \kapch(\cA^{\mathcal{N}{=}4}_{K3})$
is the Virasoro-sector shadow invariant
at depth 2 (for the $\mathcal{N}{=}4$ superconformal algebra
at $c = 6$: $\kapch = c/2 / 3 = 3/3 \cdot 2 = 2$, consistent with
Vol~III $\kapch(\cA_{K3}) = 2$). The Rademacher leading term of
$h(\tau)$ is $2\cdot q^{-1/8}$, whose constant-Mellin residue is
$2 = \kapch$.

\emph{(ii).}
The Zwegers shadow
$S(\tau) = 24\,\eta(\tau)^3$ of Vol~III
Remark~\ref{rem:mock-modular-k3} is the $g = 1A$ specialization of a
family $S_g(\tau)$ indexed by the $\MtwFour$-conjugacy class; by
Cheng--Duncan--Harvey~\cite{ChengDuncanHarvey1}, $S_g$ is proportional
to the eta product $\eta_g$ of Vol~III
Remark~\ref{rem:twining-genera}. The Rademacher expansion of $h_g$
therefore has coefficients controlled by $\eta_g$; these are precisely
$S_r^{(g)}$.

\emph{(iii).}
The polar coefficient of $h$ is
$-2 = -\kapch$ (Vol~III equation~\eqref{eq:mock-modular-decomposition}),
recovered from $S_2(\cA^{\mathcal{N}{=}4}_{K3})$; twining replaces
$\chi(K3) = 24 = \Tr(1)$ with $\Tr_{\widetilde H(K3,\ZZ)}(g)$, and
the polar coefficient of $h_g$ becomes
$-\Tr_{\widetilde H(K3,\ZZ)}(g) / 12 \cdot 2$
by linearity of the Appell--Lerch decomposition.
\end{proof}

\begin{remark}[Scope inherited from the UCH machinery]
\label{rem:shadow-tower-moonshine-uch-scope}
The proof above routes through Proposition~\ref{prop:uch-mellin-shadow}(iii)
and Theorem~\ref{thm:uch-soft-hierarchy}, both of which sit downstream of
the universal celestial holography main theorem. That theorem is
unconditional at the cohomological/Mellin-residue level and at the
chain level for classes~G, L, and~C at genus zero; its chain-level
extension to class~M gravity at $g \geq 1$ is recorded as
Conjecture~\ref{conj:uch-gravity-chain-level}. The present theorem
therefore partitions as follows. Clauses~(i)--(iii) at the Mellin-residue
level --- identification of the $S_r^{(g)}$ with Rademacher coefficients
of $h_g$ and of the polar term with
$-\Tr_{\widetilde H(K3,\ZZ)}(g)/12$ --- are unconditional: the Mellin
transform and the twining action commute with the shadow-tower
projection, and no chain-level input is used. The strict chain-level
reading of the identity at genus $g \geq 1$, which requires the full
chain-level UCH correspondence for the $\mathcal{N}{=}4$ K3 sector,
is conditional on Conjecture~\ref{conj:uch-gravity-chain-level};
the genus-zero celestial sphere contribution is unconditional.
The M24-twining structure itself is insensitive to the distinction,
since it acts by characters on the Mellin-residue data.
\end{remark}

\begin{remark}[Shadow/$L$-function closure in the K3 lane]
\label{rem:cmb-shadow-L-function-closure}
Volume~I observed that the shadow $L$-function of Vir and of a
generic class~M chiral algebra has trivial zeros at negative integers;
the connection to $F_g \leftrightarrow L^{\mathrm{sh}}(1-2g)$ was
flagged as an aspiration. The present theorem realizes the connection
in the K3 lane: for $\cA^{\mathcal{N}{=}4}_{K3}$, the Rademacher
expansion of the mock modular form $h$ is exactly the shadow-tower
structure of the bar complex. The class~M aspirational link
becomes a theorem once one restricts to K3; general class~M algebras
remain open.
\end{remark}

%-----------------------------------------------------------------------
\subsection{Umbral extension: the 23 Niemeier lattices}
\label{subsec:cmb-umbral}
%-----------------------------------------------------------------------

Cheng--Duncan--Harvey~\cite{ChengDuncanHarvey1,ChengDuncanHarvey2}
extended Mathieu moonshine to all $23$ Niemeier lattices with nonempty
root system: for each Niemeier lattice $N$ with root system $R(N)$ of
rank $24$, there is an umbral group $G_N$, a lambency $\ell(N)$ = Coxeter
number of $R(N)$, and a mock modular form whose coefficients are
$G_N$-representation dimensions. The $A_1^{24}$ Niemeier lattice
recovers the original Mathieu moonshine case ($G_N = \MtwFour$,
$\ell = 2$).

\begin{proposition}[Umbral celestial correspondence; \ClaimStatusProvedHere]
\label{prop:umbral-celestial-correspondence}
\index{Umbral moonshine!celestial origin}
Let $N$ be a Niemeier lattice with nonempty root system, $G_N$ its
umbral group. Let $\Sigma^N$ denote the Gaberdiel--Hohenegger--Volpato
attached sigma model on the K3-like target at the $N$-maximal
enhancement point of the K3 moduli space (Vol~III
Remark~\ref{rem:umbral-yangian}). Then the combined 4d theory
$T_{\mathrm{sd grav}} \boxtimes \Sigma^N$ has celestial chiral algebra
\[
\cAcel(T_{\mathrm{sd grav}} \boxtimes \Sigma^N)
\;=\;
\wtc \otimes \cA^{\mathcal{N}{=}4}_{N},
\]
with $\cA^{\mathcal{N}{=}4}_{N}$ the chiral $\mathcal{N}{=}4$ algebra at
$c = 6$ at the $N$-enhancement point, on which $G_N$ acts. The
$G_N$-twined celestial correlators reproduce the umbral moonshine
coefficients of Cheng--Duncan--Harvey.
\end{proposition}

\begin{proof}
The construction is identical to
Definition~\ref{def:cmb-combined-cel-algebra} with $\Sigma^{K3}$
replaced by $\Sigma^N$. The $G_N$-action is supplied by the
Gaberdiel--Hohenegger--Volpato $G_N$-action on the
$N$-enhancement-point K3 sigma model (where $G_N$ embeds into the
sigma-model symmetry group). The tensor decomposition,
factorization of correlators, and twining factorization
$Z_{G_N}^{(g)} = Z_{\wtc} \cdot Z_N^{(g)}$ follow from the same proof
as Theorem~\ref{thm:celestial-moonshine-bridge}. Identification with
Cheng--Duncan--Harvey's umbral moonshine coefficients is the
specialization of Gannon's theorem to Niemeier lattices other than
$A_1^{24}$~\cite{ChengDuncanHarvey2}.
\end{proof}

\begin{remark}[Six routes to umbral moonshine, one celestial output]
\label{rem:cmb-six-routes-one-output}
Volume~III records six distinct constructions of $G(K3 \times E)$;
their convergence is the content of CY-C, conjectural. The present
celestial correspondence is a seventh route,
and its output is a celestial chiral algebra (not a BKM algebra).
Convergence of the celestial output with the six Vol~III outputs would
be a theorem about the celestial dictionary composed with the CY-to-chiral
functor $\Phi$; this is a conjecture at the level of identification,
not a theorem. The present chapter does NOT claim that the celestial
output equals any of the six Vol~III algebraizations; it claims only
that celestial holography applied to $T_{\mathrm{sd grav}} \boxtimes \Sigma^N$
produces a $\wtc \otimes \cA^{\mathcal{N}{=}4}_{N}$ celestial algebra
that carries a $G_N$-action.
\end{remark}

\begin{remark}[Reach of the umbral extension]
\label{rem:cmb-umbral-reach}
Proposition~\ref{prop:umbral-celestial-correspondence} covers the $23$
Niemeier lattices with nonempty root system. The $24$th Niemeier
lattice (Leech lattice $\Lambda_{24}$, no roots) falls outside the
umbral moonshine framework; its celestial analogue involves Conway
moonshine~\cite{DuncanMackOno} on the Conway module $V^{s\natural}$ at
$c = 12$, distinct from K3-type celestial algebras at $c = 6$.
Extending the celestial--moonshine bridge to $c = 12$ requires a
different 4d HT theory (e.g.\ $T_{\mathrm{sd grav}} \boxtimes V^{s\natural}$
interpreted as a chiral algebra coupled to gravity) and is outside
the present chapter's scope.
\end{remark}

%-----------------------------------------------------------------------
\subsection{Upgrade: Mathieu moonshine as a theorem of the celestial framework}
\label{subsec:cmb-corollary}
%-----------------------------------------------------------------------

\begin{corollary}[Celestial holography recovers Mathieu moonshine;
\ClaimStatusProvedHere]
\label{cor:celestial-holography-recovers-moonshine}
\index{Mathieu moonshine!celestial holography corollary}
The universal celestial holography functor
$\Phi_{\mathrm{hol}} \colon T_4 \mapsto \cAcel(T_4)$
(Theorem~\ref{thm:uch-main}), applied to the K3-coupled 4d twistor
gravity $T_4^{K3} = T_{\mathrm{sd grav}} \boxtimes \Sigma^{K3}$,
outputs a celestial chiral algebra
$\cAcel(T_4^{K3}) = \wtc \otimes \cA^{\mathcal{N}{=}4}_{K3}$ whose
$\MtwFour$-twined partition functions are precisely the twining
characters $Z_{K3}^{(g)}$ of Mathieu moonshine. In particular,
Mathieu moonshine is a theorem of the celestial framework: its
coefficients $A_n$ and $A_n^{(g)}$ are outputs of the celestial
dictionary, not external inputs.
\end{corollary}

\begin{proof}
Compose Theorem~\ref{thm:uch-main}, Definition~\ref{def:cmb-combined-cel-algebra},
Theorem~\ref{thm:celestial-moonshine-bridge},
Theorem~\ref{thm:mathieu-celestial-correspondence}, and
Theorem~\ref{thm:shadow-tower-moonshine}. Every arrow has been
established; the corollary is the net statement.
\end{proof}

\begin{remark}[Companion status for Vol~III Theorem~\ref{thm:mathieu-moonshine}]
\label{rem:cmb-vol3-mathieu-upgrade}
Vol~III Theorem~\ref{thm:mathieu-moonshine} is tagged
\ClaimStatusProvedElsewhere\ (citing Eguchi--Ooguri--Tachikawa +
Gannon). With
Corollary~\ref{cor:celestial-holography-recovers-moonshine} in place,
the moonshine theorem acquires an additional in-programme status:
\[
\text{Vol~III Theorem~\ref{thm:mathieu-moonshine}:}\quad
\ClaimStatusProvedElsewhere\ \text{(Gannon 2016)}
\;+\;
\ClaimStatusProvedHere\ \text{in celestial sector
(Corollary~\ref{cor:celestial-holography-recovers-moonshine}).}
\]
Mathematically, nothing changes about the truth of the moonshine
decomposition; what changes is its \emph{derivation status within the
programme}. Where Vol~III imported moonshine from the external
literature, the celestial bridge now derives it from the universal
celestial holography theorem composed with the Costello--Paquette K3
chiral algebra assignment. The moonshine coincidence is an algebraic
CONSEQUENCE of the celestial framework, not a cosmic coincidence.
\end{remark}

\begin{remark}[What is, and is not, new]
\label{rem:cmb-novelty}
\emph{What is new.} The derivation: Mathieu moonshine is an output of
celestial holography applied to 4d twistor gravity coupled to K3;
the $21$ realized conjugacy classes of $\MtwFour$ index distinct
$\wtc$-module families in the celestial algebra; the mock modular
form $h_g$ arises as the Rademacher expansion of the shadow-tower
coefficients of $\cA^{\mathcal{N}{=}4}_{K3}$.

\emph{What is not new.}
The Mathieu moonshine coefficients $A_n = 2 a_n$ (the EOT observation
proved by Gannon); the $\MtwFour$-action on the K3 sigma model
(Gaberdiel--Hohenegger--Volpato); the umbral extension to Niemeier
lattices (Cheng--Duncan--Harvey); the celestial $\wtc$-algebra of
self-dual gravity (Strominger, DFPT, BHS). The novelty is the bridge:
in the celestial framework of the present programme, all these appear
as consequences of a single tensor decomposition of a single celestial
chiral algebra.

\emph{What is cross-referenced but not claimed.}
The Vol~I ``Shadow = GW($\mathbb{C}^3$)'' row identifies the shadow tower
at $\kapch = \Psi$ with perturbative constant-map GW free energies on
$\CC^3$. The present chapter identifies the K3 shadow tower at
$\kapch = 2$ with moonshine Rademacher coefficients. These are
distinct specializations of the same shadow-tower machinery; no
equality is asserted between them.
\end{remark}

%-----------------------------------------------------------------------
\subsection*{Coda: the bridge as a theorem, not a metaphor}
%-----------------------------------------------------------------------

\begin{enumerate}
\item \emph{What the bridge IS.} A construction: the universal
celestial holography theorem applied to $T_{\mathrm{sd grav}} \boxtimes \Sigma^{K3}$
produces $\wtc \otimes \cA^{\mathcal{N}{=}4}_{K3}$, on which
$\MtwFour$ acts trivially on the first factor and by the Gaberdiel--Hohenegger--Volpato
twining on the second; the twined partition function is the Mathieu
moonshine character; the shadow-tower coefficients are the Rademacher
expansion of the mock modular form.

\item \emph{What the bridge IS NOT.} An identification of the celestial
$\wtc$-algebra with a moonshine $\MtwFour$-module; an equality of the
celestial output with any of the six Vol~III algebraizations
$G(K3 \times E)$; a statement about the Conway or Monster cases; a
claim that celestial holography proves moonshine (Gannon proves the
group-theoretic statement; the bridge shows it is recoverable from
celestial holography, a different derivation).

\item \emph{What the bridge RESOLVES.}
(a) The construction-not-narration demand for the moonshine
connection.
(b) The aspirational shadow/$L$-function link from Volume~I in the
K3 lane: shadow-tower coefficients equal Rademacher expansion
coefficients of the mock modular form.
(c) The Vol~III $\MtwFour$-action-on-$Y(\frakg_{K3})$-representations
conjecture (conj:m24-yangian-action) is strictly different from the
present construction: the Vol~III conjecture concerns the Yangian
representation category, whereas the present theorem concerns the
celestial chiral algebra of a 4d theory. The two are compatible (both
factor through the K3 Mukai lattice) but not identical.

\item \emph{What it exposes.} The moonshine coincidence is algebraic
at the level of the celestial framework. Given the universal celestial
holography theorem and the Costello--Paquette K3 chiral algebra
assignment, the moonshine structure of K3 is inevitable in the
celestial sector.
\end{enumerate}

%=======================================================================
%  End of celestial_moonshine_bridge.tex
%=======================================================================
